\chapter{The owner policy}
\label{ch:policy}

The policy is mandatory on every worker. Its trigger is local \textbf{work}
appearing --- not a human being detected. When the machine's own workload grows,
Condor jobs suspend; if contention persists, they are evicted.

\section{The two tiers}

\begin{lstlisting}
NonCondorLoad = (TotalLoadAvg - TotalCondorLoadAvg)

# TIER 1 - admission and suspension (costs nothing, loses nothing)
START = (NonCondorLoad   < TotalCpus * 0.50)      # is the OWNER busy?
     && (TotalLoadAvg    < TotalCpus * 0.90)      # is the MACHINE busy?
     && (HostCpuPsiAvg10 < 20.0)                  # actual CPU stall, 10s window
     && (HostMemAvailMB  > HostMemTotalMB  * 0.25)
     && (HostDiskAvailMB > HostDiskTotalMB * 0.10)
     && FITS_CPU && FITS_MEM && FITS_DISK         # does THIS job fit?

FITS_MEM = (TARGET.RequestMemory =?= UNDEFINED)
        || (TARGET.RequestMemory < (HostMemAvailMB - MEM_SOFT_FLOOR))

WANT_SUSPEND = !(EVICT_PRESSURE)     # selects the mechanism - see chapter 5
SUSPEND      = (NonCondorLoad > TotalCpus * 0.75)
CONTINUE     = (NonCondorLoad < TotalCpus * 0.50)

# TIER 2 - eviction (the only thing that discards work)
PREEMPT = EVICT_PRESSURE =
          (HostMemAvailMB > 0)
       && (  (HostMemAvailMB  < max(1GB, HostMemTotalMB * 0.10))
          || (HostMemPsiAvg10 > 10.0)
          || (HostDiskAvailMB < HostDiskTotalMB * 0.05) )

MaxJobRetirementTime = 0
MachineMaxVacateTime = 30
\end{lstlisting}

\section{Admission is gated on measurement, never on accounting}

This is the central design decision of the whole system.

\begin{keypoint}
Condor's slot accounting is \textbf{request}-based. A job asking for one core may
spawn eight threads; a job asking for 256\,MB may use 10\,GB. Sixteen
``one-core'' jobs can saturate a sixteen-core box while matchmaking still
believes the machine is empty.
\end{keypoint}

Every term in the policy above comes from \file{/proc}. Nothing in it trusts what
a job claimed it would use, because that number is an aspiration, not a
constraint.

\subsection{The two load terms are not redundant}

\kn{NonCondorLoad} subtracts Condor's own contribution and answers \emph{is the
owner busy}. \kn{TotalLoadAvg} includes it and answers \emph{is the machine
busy}. Both are needed.

With only the first, a box at load 60 driven entirely by Condor reads as idle
and keeps admitting more work: the owner stays perfectly protected while the
machine drowns. \S\ref{sec:closure} shows this measured.

\subsection{Admission is graduated, not binary}

Comparing each job's request against measured headroom lets a small job land
where a large one cannot. This is what the \kn{FITS\_*} terms do.

True dynamic advertisement is not available. \kn{Cpus} and \kn{Memory} on a
partitionable slot are configured capacity minus request-based allocations,
computed by the startd itself, and cannot safely be overridden. So
\kn{condor\_status} shows nominal capacity while \kn{START} holds the truth.

\begin{trap}
\textbf{A machine can display free slots and still decline every job.} This is
not a malfunction, it is the design --- but it means \kn{condor\_status} alone
cannot tell you whether a machine is working. Diagnose with:

\begin{lstlisting}[style=term]
condor_status -af Name Start TotalLoadAvg HostCpuPsiAvg10 \
                       HostMemAvailMB HostDiskAvailMB
\end{lstlisting}
\end{trap}

\section{Pressure Stall Information}

\kn{PSI} is the signal that means \emph{the owner is suffering}.

Free memory and load average are proxies. \file{/proc/pressure/*} measures the
percentage of time tasks actually stalled waiting for a resource. A box can
report gigabytes available while thrashing, and load average is a decaying
one-minute mean that lets a negotiation cycle admit a whole burst before it
registers. PSI \kn{avg10} reacts in ten seconds.

\section{Thresholds scale with the machine}

The fleet spans a 4\,GB single-board computer and a 31\,GB workstation. A single
fraction is wrong at one end or the other:

\begin{itemize}
  \item 10\% of a 4\,GB SBC is 410\,MB --- too little to avoid thrashing.
  \item 25\% of a 30\,GB box is 7.5\,GB --- needlessly idle.
\end{itemize}

\kn{ifThenElse} gives one expression that is correct at both ends:

\begin{lstlisting}
MEM_HARD_FLOOR = ifThenElse((HostMemTotalMB * 0.10) > 1024,
                            (HostMemTotalMB * 0.10), 1024)
\end{lstlisting}

\section{CPU contention suspends. Memory pressure evicts. Time does nothing.}

\begin{keypoint}
A suspended job costs the owner \textbf{no CPU at all} --- \kn{SIGSTOP} means it
is not scheduled --- so discarding its work merely because time has passed is
pure waste.

Memory is the one resource suspension cannot return: \kn{SIGSTOP} freezes a
process but keeps its RSS. Memory pressure is therefore the \emph{only}
justification for killing.
\end{keypoint}

When eviction does fire, reclaim is immediate: \kn{MaxJobRetirementTime = 0},
then 30 seconds to exit cleanly before \kn{SIGKILL}.

Thresholds are busy cores measured against the reservation, not arbitrary
numbers. On a four-core box donating two, \kn{SUSPEND} at $>2.0$ means ``the
owner now needs more than the two cores reserved for them''. \kn{CONTINUE} at
$<1.0$ is the hysteresis that stops jobs flapping.

\section{No per-job memory ceiling}

\kn{CGROUP\_MEMORY\_LIMIT\_POLICY} is set to \kn{none}, deliberately.

The alternative, \kn{hard}, enforces \kn{request\_memory} as an absolute per-job
limit --- so a job asking for 256\,MB is OOM-killed at 256\,MB with 14\,GB free.
On machines of 4 to 16\,GB that forces everyone to over-request, wasting exactly
the RAM the fleet does not have, and turns honest requests into spurious
failures.

\kn{request\_memory} therefore stays a \textbf{matchmaking figure only}. Memory
is protected system-wide by eviction instead.

The margin is a fraction of installed RAM rather than an absolute floor: 2\,GB is
12\% of a 16\,GB desktop and half of a 4\,GB SBC, and this fleet spans both.

\kn{HostMemAvailMB} and \kn{HostMemTotalMB} come from the \kn{STARTD\_CRON} job,
because Condor publishes no free-memory attribute of its own. \kn{MemAvailable}
is used rather than \kn{MemFree}, which ignores reclaimable page cache and would
report false pressure on any long-running machine.

\subsection{Two traps in the eviction expression}

\begin{trap}
\textbf{Do not gate eviction on \kn{Activity == "Suspended"}.} \kn{SUSPEND} fires
on CPU load, so a job quietly eating RAM without loading the CPU would never
suspend and therefore never be preempted --- unenforceable in exactly the case
that matters.
\end{trap}

\begin{trap}
\textbf{Guard with \kn{HostMemAvailMB > 0}.} Before the cron's first run the
attribute is 0 or \kn{UNDEFINED}, and a bare comparison reads as pressure ---
evicting every job the moment the startd starts.
\end{trap}

\section{Presence must not gate admission}

\kn{OwnerSessionActive} is published and useful for diagnostics, but it is
deliberately \textbf{not} in \kn{START}.

People do not log out. On any desktop the attribute reads true essentially
forever, so gating admission on it makes the machine advertise its cores and
refuse everything, permanently --- and a machine refusing everything is
indistinguishable, in \kn{condor\_status}, from a healthy idle one.

The policy reacts to local \emph{work}, not to a human existing. If the owner
starts doing something, load and PSI catch it within seconds. Their mere
presence costs them nothing, because Condor jobs already yield the CPU instantly.

\begin{keypoint}
This only surfaced when the first machine that is both \textbf{small} and
\textbf{actually in use} joined the pool. Servers and headless boxes never expose
it --- which is why a pool tested only on spare hardware will ship this bug.
\end{keypoint}

\section{The policy is inert without the stanza that feeds it}
\label{sec:inert}

Every attribute the policy references comes from the \kn{STARTD\_CRON} job.

\begin{trap}
Without that stanza the configuration still \textbf{parses}; \kn{condor\_config\_val}
still reports the correct expressions; \kn{condor\_status} still shows a healthy
\st{Unclaimed/Idle} slot --- while \kn{START} evaluates to \kn{UNDEFINED} and both
admission control and eviction are silently dead.
\end{trap}

\kn{UNDEFINED} is \textbf{worse than false}. A false \kn{START} parks the slot in
\st{Owner} state, which is visible. Undefined leaves it displaying
\st{Unclaimed/Idle}, indistinguishable from healthy.

This was introduced \emph{twice}, by edits that replaced the policy section
wholesale and took the stanza with it. Two mitigations follow from that: keep the
stanza \textbf{above} the policy header, and after any configuration edit verify
the attributes are on the \textbf{slot ad}, not merely in the configuration:

\begin{lstlisting}[style=term]
condor_status -af Name Start OwnerSessionActive HostMemAvailMB HostDiskAvailMB
\end{lstlisting}

Note that the collector only sees a changed ad on the startd's next update, so a
newly published attribute can read \kn{undefined} for several minutes without
anything being wrong. \kn{condor\_reconfig} forces a fresh ad.

\section{Do not write this policy against ConsoleIdle or KeyboardIdle}
\label{sec:console}

On a \textbf{partitionable} slot, machine-level attributes behave differently on
the dynamic slots where jobs actually run:

\begin{table}[H]
\centering
\begin{tabularx}{\textwidth}{L{0.34\textwidth} L{0.22\textwidth} Y}
\toprule
\textbf{Attribute} & \textbf{Partitionable parent} & \textbf{Dynamic slot} \\
\midrule
\kn{ConsoleIdle} & real value & \textbf{absent (UNDEFINED)} \\
\kn{KeyboardIdle} & real value & \kn{-1} sentinel \\
\kn{LoadAvg} & real value & \kn{-1.0} sentinel \\
\kn{TotalLoadAvg}, \kn{TotalCondorLoadAvg} & real & \textbf{real} \\
\kn{STARTD\_CRON} attributes & present & \textbf{present} \\
\bottomrule
\end{tabularx}
\caption{Attribute availability differs between the partitionable parent and its dynamic slots.}
\end{table}

One \kn{UNDEFINED} term makes an entire \kn{\&\&} chain non-true. A policy written
against \kn{ConsoleIdle} therefore leaves every dynamic slot in state \st{Owner},
refusing every claim --- while the parent still reports \st{Unclaimed/Idle} with
\kn{Start = true}.

\begin{trap}
\textbf{This failure is invisible from \kn{condor\_status}.} The pool looks
healthy, jobs sit \st{Idle} forever, and \kn{condor\_q -better-analyze} reports
``1 are able to run your job''. Only \file{StartLog} shows
\kn{Request to claim resource refused} and \kn{Claiming protocol failed}.
\end{trap}

Use the machine-wide \kn{Total*} attributes instead.

\kn{KeyboardIdle} is separately wrong on any administered machine: it counts
activity on \emph{any} tty, including ssh ptys, so an admin session makes the box
look permanently busy. Measured with nobody present: \kn{KeyboardIdle = 2} while
\kn{ConsoleIdle = 359657}.

\section{Detecting remote users}

\kn{condor\_kbdd} only sees the local X console, so on machines used over RDP or
VNC it is blind to the actual user.

A \kn{STARTD\_CRON} job publishing \kn{OwnerSessionActive} --- derived from
established connections on the remote-desktop port plus active non-greeter
graphical sessions --- covers that case. Unlike the console attributes, it
\emph{does} propagate to dynamic slots.

\begin{table}[H]
\centering
\begin{tabularx}{\textwidth}{L{0.34\textwidth} Y}
\toprule
\textbf{Setting} & \textbf{Rationale} \\
\midrule
\kn{CGROUP\_MEMORY\_LIMIT\_POLICY = none} & No per-job cap; memory is defended system-wide \\
\kn{JOB\_RENICE\_INCREMENT = 10} & Covers the lag: load average is a decaying one-minute mean and reacts in tens of seconds. Renice makes jobs yield the CPU instantly \\
\kn{MachineMaxVacateTime = 30} & Long enough to exit cleanly, short enough that the owner is not kept waiting \\
\kn{KILL = FALSE} & Never hard-kill on policy; let the vacate path run \\
\bottomrule
\end{tabularx}
\caption{Supporting settings and why each is set as it is.}
\end{table}
